\chapter{Recherche bibliographique et état de l'art}

\section{Démarche de recherche}

Cette étude ne prétend pas être une revue systématique de toute la
littérature sur l'IIoT. Elle répond à une question plus concrète : quelles
solutions permettent de construire une supervision industrielle à partir de
données déjà disponibles sur une plateforme edge ? La recherche a donc suivi
six axes : l'interopérabilité OT/IT, la surveillance conditionnelle, le
transport MQTT, l'ingestion des messages, le stockage temporel et la
visualisation.

Les sources ont été choisies dans un ordre simple. Les normes et
spécifications officielles ont servi pour les protocoles. Les documentations
des éditeurs ont permis de vérifier les possibilités des outils. Enfin, des
articles scientifiques ont été consultés pour situer le projet par rapport à
la maintenance conditionnelle et prédictive. Les critères de sélection ont
été la compatibilité avec Litmus Edge, la possibilité d'un déploiement local,
la lisibilité pour une équipe industrielle et un coût raisonnable pour une
seule ligne.

\section{De la maintenance corrective à la surveillance conditionnelle}

La maintenance corrective intervient après la panne. La maintenance
préventive remplace ou contrôle un équipement selon un calendrier, même si son
état réel ne l'exige pas encore. La maintenance conditionnelle s'appuie au
contraire sur des mesures : température, pression, vibration, niveau,
viscosité ou compteur d'heures. Elle cherche une dérive avant que le défaut ne
bloque la production.

La maintenance prédictive va plus loin. Elle utilise un historique suffisant
pour estimer une probabilité de panne ou une durée de vie restante. La revue
de Krupitzer et al. montre que cette démarche repose sur la collecte, la
préparation et l'interprétation de données industrielles, souvent avec des
méthodes d'apprentissage automatique~\cite{krupitzer2020survey}. Li, Fei et
Zhang décrivent une architecture edge-to-cloud appliquée à la surveillance
d'équipements de fabrication. Le traitement à la périphérie réduit le volume
échangé et rapproche l'analyse de la machine~\cite{li2022condition}.

GlassInk Monitor se place volontairement au niveau de la \textbf{surveillance
conditionnelle}. Le système détecte des seuils, suit des tendances et
extrapole le temps avant épuisement d'un fluide. Il ne prédit pas encore une
panne par apprentissage automatique. Il manque pour cela des mois de données
réelles, des défauts correctement étiquetés et un retour de la maintenance.
Employer dès maintenant le terme de maintenance prédictive donnerait une
image plus ambitieuse, mais techniquement inexacte.

\section{Données disponibles sur la Markem-Imaje 9450c}

La ligne étudiée utilise une imprimante Markem-Imaje 9450c à jet d'encre
continu, adaptée aux cadences industrielles et équipée d'une connectivité
Ethernet~\cite{markem9450}. Elle est raccordée à un automate Siemens
S7-1515-2~PN. La communication s'appuie sur TCP/IP et sur des trames contenant
un identifiant, une longueur, des données utiles et une somme de contrôle.

\begin{figure}[H]
  \centering
  \includegraphics[width=0.66\textwidth]{imprimante-markem-imaje-9450.png}
  \caption{Imprimante CIJ Markem-Imaje 9450 et tête de marquage}
  \label{fig:imprimante-9450}
\end{figure}

Le périmètre de supervision retient huit familles d'informations : alertes,
état général, état de l'impression, état du jet, informations des deux
cartouches, paramètres de fonctionnement et tâche d'impression active.

La réponse des paramètres de fonctionnement contient 84 octets et plus de
quarante valeurs. Parmi elles figurent la vitesse moteur, la pression, la
viscosité, le niveau d'encre, la quantité d'additif injectée, les températures
électronique, encre et tête, les états des électrovannes, le niveau du tube,
l'autonomie d'encre et la consommation moyenne. Les cartouches externes
d'encre et d'additif ont une capacité de 800~mL ; les réservoirs internes ont
une capacité annoncée de 1,075~L et maintiennent une réserve de fonctionnement
d'environ 24~heures. La capacité des cartouches externes et les fonctions de
suivi de l'autonomie sont également indiquées dans la documentation technique
du constructeur~\cite{markem9450}.

Dans l'architecture existante, la 9450c est reliée à Logoface4, puis au système
Gosun. Les données rejoignent ensuite Litmus Edge et l'IoT Backbone. Cette
organisation confirme que la plateforme de
supervision n'a pas à interroger directement l'imprimante : les données déjà
centralisées sur Litmus Edge constituent son point d'entrée. Elle fournit aussi
une base concrète au simulateur, qui reproduit les codes d'état du jet, l'état
d'impression et les paramètres utiles à la maintenance.

\section{Acquisition et interopérabilité à la frontière OT/IT}

Dans une architecture industrielle classique, OPC UA offre un modèle
d'information commun, des services client--serveur et un mode
publication--abonnement. La spécification prévoit aussi l'authentification,
le chiffrement et le contrôle d'intégrité~\cite{opcua}. Une application
pourrait donc lire directement les variables d'un automate ou d'un serveur
OPC UA.

Ce choix n'est pas retenu ici. Les automates, l'imprimante et les équipements
de la ligne sont déjà raccordés à Litmus Edge. La plateforme communique avec
les équipements, conserve des données localement et peut les publier vers
MQTT ou REST~\cite{litmusedge}. Ajouter un second client direct vers les PLC
dupliquerait la collecte et agrandirait inutilement le périmètre de sécurité.
Le point d'entrée de GlassInk Monitor est donc la sortie nord de Litmus Edge.
L'application reste en lecture seule et ne commande aucun équipement.

\section{Transport des données}

MQTT découple la source et les consommateurs. Litmus Edge publie sans connaître
Node-RED, et Node-RED peut redémarrer sans modifier la configuration des
automates. Le QoS~1 retenu garantit une livraison au moins une fois, ce qui
signifie qu'un message peut être reçu deux fois~\cite{mqtt5}. Cette propriété
explique la protection contre les doublons mise en place dans la base.

MQTT ne définit cependant ni le nom des topics industriels ni la structure du
contenu. Sparkplug complète MQTT avec un espace de noms, un format de charge
utile et une gestion de l'état de session adaptée aux usages
IIoT~\cite{sparkplug}. Son format binaire Protobuf demande néanmoins un
décodage supplémentaire.

Pour le prototype, des topics explicites et un contenu JSON ont été choisis.
Ils sont faciles à lire, à simuler et à contrôler pendant une démonstration.
Lors du raccordement, deux cas restent possibles : Litmus Edge publie ce JSON,
ou il publie en Sparkplug B et un décodeur est ajouté à l'entrée du flux.
Cette adaptation ne change ni le modèle des mesures ni les tableaux de bord.

\section{Solutions d'ingestion comparées}

L'ingestion ne consiste pas seulement à copier un message. Elle vérifie le
schéma, sépare télémétrie et événements, rejette les données invalides et
garantit l'idempotence. Trois approches ont été étudiées.

\begin{table}[H]
\centering
\small
\caption{Comparaison des solutions d'ingestion}
\begin{tabularx}{\textwidth}{>{\raggedright\arraybackslash}p{0.17\textwidth}YYY}
\toprule
Solution & Points forts & Limites dans ce projet & Appréciation \\
\midrule
Telegraf &
Configuration déclarative, nombreux plugins, faible consommation. &
Moins souple pour les événements hétérogènes, les remplissages et les
transactions métier. &
Très adapté à des mesures régulières et déjà normalisées. \\
Node-RED &
Flux visuel, adaptation rapide des messages, bonne lisibilité du chemin de
données. &
L'éditeur doit être protégé et les flux doivent rester versionnés et testés. &
Retenu pour le prototype et le raccordement à Litmus Edge. \\
Service spécifique &
Contrôle complet du décodage et des transactions. &
Davantage de code, de maintenance et une logique moins visible pour l'équipe
industrielle. &
À réserver à un besoin que Node-RED ne peut pas couvrir proprement. \\
\bottomrule
\end{tabularx}
\end{table}

Telegraf reste une bonne solution lorsque les messages sont déjà homogènes et
que le chemin entre la source et la base demande peu d'adaptation~\cite{telegraf}.
Ici, la télémétrie, les événements de maintenance et les résultats de marquage
n'ont pas la même structure. Node-RED a donc été retenu pour valider chaque
famille, ajouter les tags utiles et produire le Line Protocol attendu par
InfluxDB. Le flux reste lisible dans l'éditeur tout en étant conservé dans le
dépôt du projet.

\section{Stockage, visualisation et alertes}

InfluxDB organise une donnée autour d'une mesure, de tags indexés, de champs et
d'un horodatage~\cite{influxdb,influxschema}. Ce modèle correspond directement
aux variables de l'imprimante : l'identité de la machine et la ligne deviennent
des tags, tandis que les niveaux, températures et compteurs restent des champs.
Les défauts et les marquages sont conservés dans des mesures distinctes. Cette
séparation évite de multiplier les colonnes lorsque la liste des tags Litmus
évolue.

Le prototype utilise InfluxDB OSS 2.7 avec une organisation et un bucket dédiés.
La durée de rétention du bucket est illimitée pour la démonstration ; elle devra
être fixée avec les équipes qualité et infrastructure avant le raccordement au
site. Node-RED enrichit chaque point avec la classification de la machine et les
prévisions de consommables. Les requêtes Flux filtrent ensuite les séries,
appliquent les fenêtres temporelles et agrègent les indicateurs affichés dans
Grafana~\cite{flux}.

Grafana dispose d'une source InfluxDB native. Elle interroge le bucket en Flux,
construit les vues opérateur, maintenance et direction, puis évalue les règles
d'alerte~\cite{grafanainflux}. Pour une ligne et moins de 200 pièces par heure,
une instance de chaque service suffit. La haute disponibilité pourra être
étudiée à partir de mesures réelles de charge et des exigences du site.

\section{Synthèse et positionnement du projet}

\begin{table}[H]
\centering
\small
\caption{Positionnement de GlassInk Monitor}
\begin{tabularx}{\textwidth}{
  >{\raggedright\arraybackslash}p{0.16\textwidth}
  >{\raggedright\arraybackslash}p{0.25\textwidth}
  >{\raggedright\arraybackslash}p{0.20\textwidth}
  Y}
\toprule
Brique & Options étudiées & Choix & Motif principal \\
\midrule
Point d'intégration &
PLC/OPC UA direct, Litmus Edge &
Litmus Edge &
La collecte industrielle existe déjà. \\
Transport &
MQTT JSON, Sparkplug B &
MQTT JSON au prototype &
Lecture et simulation simples ; décodage Sparkplug ajoutable. \\
Ingestion &
Telegraf, Node-RED, service spécifique &
Node-RED &
Transformation visuelle des mesures et événements. \\
Stockage &
InfluxDB, base relationnelle, stockage par fichiers &
InfluxDB 2.7 &
Modèle temporel natif, rétention et intégration directe avec Grafana. \\
Supervision &
Application web, Grafana &
Grafana &
Tableaux de bord, rôles et alertes déjà disponibles. \\
Déploiement &
Services natifs, Docker Compose, Kubernetes &
Docker Compose &
Reproductible et proportionné à une seule ligne. \\
Niveau d'analyse &
Seuils, surveillance conditionnelle, modèle prédictif &
Surveillance conditionnelle &
Pas encore d'historique réel étiqueté pour entraîner un modèle. \\
\bottomrule
\end{tabularx}
\end{table}

L'état de l'art conduit ainsi à une architecture edge-to-dashboard maîtrisée :
Litmus Edge comme frontière industrielle, MQTT pour le découplage, Node-RED
pour l'adaptation, InfluxDB pour l'historisation et Grafana pour l'usage
quotidien. L'apport du projet tient à la manière dont ces briques répondent au
fonctionnement de la ligne : l'activité dépend de la commande, les fluides
peuvent expirer avant d'être vides et le raccordement final doit rester limité
à la configuration de la source Litmus Edge.

\clearpage
\section{Problématique retenue}

À la suite de cette analyse, la problématique du PFA est formulée ainsi :

\begin{projectbox}
\centering
\emph{Comment transformer les données disponibles sur Litmus Edge en une
supervision fiable de l'imprimante CIJ, capable d'anticiper les ruptures et la
péremption des consommables, de distinguer un arrêt réel d'une période sans
commande et de présenter une information adaptée à chaque utilisateur ?}
\end{projectbox}

Cette question générale se décline en trois questions d'ingénierie :

\begin{enumerate}
  \item comment transporter et enregistrer les messages en tenant compte des
        retransmissions possibles avec le QoS~1 ?
  \item comment calculer la disponibilité sans classer une absence de commande
        comme une panne ?
  \item comment réunir dans une seule décision le temps avant épuisement et la
        date effective de péremption d'un fluide ?
\end{enumerate}

Deux contraintes encadrent la réponse. Le projet ne commande ni l'imprimante
ni les automates : le flux reste en lecture seule. Par ailleurs, l'accès au
Litmus Edge réel n'était pas disponible pendant le développement. La solution
devait donc être construite et vérifiée avec un simulateur, puis rester
remplaçable par la source industrielle sans reprendre toute l'application.
